\documentclass[11pt]{article}

% Sample final exam (public): Lectures 3 and 4 as taught (Sections 3.1-3.12; 4.1-4.7, 4.9 up to confidence intervals, 4.12-4.14), using Lectures 1-2
\usepackage[margin=1in]{geometry}
\usepackage[T1]{fontenc}
\usepackage{lmodern}
\usepackage{amsmath,amssymb,amsthm}
\usepackage{enumitem}
\usepackage{hyperref}
\usepackage{xcolor}

\title{Sample Final Exam}
\author{Lectures 3--4: Vector Calculus, Optimization, Probability and Statistics}
\date{}

\theoremstyle{definition}
\newtheorem{problem}{Problem}

\newcommand{\norm}[1]{\lVert #1 \rVert}
\newcommand{\ip}[2]{\langle #1, #2 \rangle}
\newcommand{\pts}[1]{\textup{[#1~pts]}}
\newcommand{\E}{\mathbb{E}}

\DeclareMathOperator{\tr}{tr}
\DeclareMathOperator{\diag}{diag}
\DeclareMathOperator{\Var}{Var}
\DeclareMathOperator{\Cov}{Cov}
\DeclareMathOperator{\Corr}{Corr}
\DeclareMathOperator{\KL}{KL}

\begin{document}

\maketitle
\thispagestyle{empty}

\noindent Name: \rule{7cm}{0.4pt} \hfill Student ID: \rule{4cm}{0.4pt}

\medskip
\noindent\fbox{\parbox{\dimexpr\linewidth-2\fboxsep-2\fboxrule}{
\textbf{Time:} 90 minutes. \textbf{Total:} 100 points in Problems 1--5. Problem 6 is a bonus problem for up to 10 extra points; the exam score is capped at 100.\\
\textbf{Allowed:} the final cheatsheet handed out with this exam. No notes, books or calculators; the numerical values you need are given in the problems.\\
Justify every answer. You may quote any result from the lectures by name or by its statement on the cheatsheet. Vectors of $\mathbb{R}^n$ are columns, $\sigma(t) = \frac{1}{1 + e^{-t}}$ is the sigmoid function, and $\log$ is the natural logarithm.
}}

\bigskip

\begin{problem}[20 points]\label{pr:fin-logistic}
    Let $\mathbf{x}_1, \ldots, \mathbf{x}_N \in \mathbb{R}^n$ be data points with labels $y_1, \ldots, y_N \in \{0, 1\}$, let $X \in \mathbb{R}^{N \times n}$ have the rows $\mathbf{x}_i^\top$, and let $\lambda \geq 0$. Regularized logistic regression minimizes
    \[J(\mathbf{w}) = \frac{1}{N}\sum_{i=1}^N \big(-y_i\log p_i - (1 - y_i)\log(1 - p_i)\big) + \frac{\lambda}{2}\norm{\mathbf{w}}^2, \qquad p_i = \sigma(\mathbf{w}^\top\mathbf{x}_i).\]
    \begin{enumerate}[label=(\alph*)]
        \item \pts{5} Show that $\nabla J(\mathbf{w}) = \frac{1}{N}X^\top(\mathbf{p} - \mathbf{y}) + \lambda\mathbf{w}$ and $\nabla^2 J(\mathbf{w}) = \frac{1}{N}X^\top DX + \lambda I$, where $\mathbf{p} = (p_1, \ldots, p_N)^\top$ and $D = \diag\big(p_1(1 - p_1), \ldots, p_N(1 - p_N)\big)$. You may use $\sigma' = \sigma(1 - \sigma)$.
        \item \pts{5} Let $\lambda > 0$. Show that every eigenvalue of every Hessian $\nabla^2 J(\mathbf{w})$ lies in the interval $[\lambda, \ L]$ with $L = \frac{\norm{X}_2^2}{4N} + \lambda$. Conclude that $J$ is $\lambda$-strongly convex and $L$-smooth, and that it has exactly one minimizer $\mathbf{w}^*$.
        \item \pts{6} Let $n = 2$, let the four data points be $(1, 1)$, $(1, -1)$, $(-1, 1)$ and $(-1, -1)$ with arbitrary labels, and let $\lambda = 0.05$. Compute the constants of part (b), and find a number $k$ of gradient descent steps with step size $\eta = \frac{1}{L}$ which guarantees $J(\mathbf{w}_k) - J(\mathbf{w}^*) \leq 10^{-6}\big(J(\mathbf{w}_0) - J(\mathbf{w}^*)\big)$ for every $\mathbf{w}_0$. (Use $1 - t \leq e^{-t}$ and $\log 10 \approx 2.303$.)
        \item \pts{4} Let $\lambda = 0$, and let the data be linearly separable: there is a $\mathbf{u}$ with $\mathbf{u}^\top\mathbf{x}_i > 0$ whenever $y_i = 1$ and $\mathbf{u}^\top\mathbf{x}_i < 0$ whenever $y_i = 0$. Show that $J$ has no minimizer.
    \end{enumerate}
\end{problem}

\begin{problem}[20 points]\label{pr:fin-matrix-calculus}
    \begin{enumerate}[label=(\alph*)]
        \item \pts{7} Let $B \in \mathbb{R}^{n \times p}$ be fixed, and let $F: \mathbb{R}^{m \times n} \to \mathbb{R}^{m \times p}$, $F(X) = XB$. Its derivative is a tensor of order $4$ with the entries $\frac{\partial F_{ij}}{\partial X_{kl}}$. Compute these entries (use the Kronecker delta). For a scalar loss $L(F(X))$ with $G = \nabla_F L \in \mathbb{R}^{m \times p}$, use the chain rule in Einstein notation to show that $\nabla_X L = GB^\top$, and write this contraction as a call of \texttt{np.einsum} that takes $G$ and the order-$4$ array \texttt{DF}.
        \item \pts{6} Let $S$ be a symmetric positive definite $d \times d$ matrix, and let
        \[\varphi(\Sigma) = \log\det\Sigma + \tr(S\Sigma^{-1})\]
        for symmetric positive definite $\Sigma$. Treating the $d^2$ entries of $\Sigma$ as independent variables, compute $\nabla\varphi(\Sigma)$, and show that $\Sigma = S$ is the only critical point.
        \item \pts{7} Show that $\Sigma = S$ is the global minimum of $\varphi$. (Hint: write $\Sigma = S^{1/2}MS^{1/2}$, and express $\varphi(\Sigma) - \varphi(S)$ through the eigenvalues of $M$.) Explain why this proves that the maximum likelihood estimate of the covariance matrix of a Gaussian, with the mean estimated by $\bar{\mathbf{x}}$, is the sample covariance matrix $S = \frac{1}{N}\sum_i(\mathbf{x}_i - \bar{\mathbf{x}})(\mathbf{x}_i - \bar{\mathbf{x}})^\top$, when $S$ is invertible.
    \end{enumerate}
\end{problem}

\begin{problem}[20 points]\label{pr:fin-gaussian}
    Let $(X, Y)^\top \sim \mathcal{N}(\mathbf{0}, \Sigma)$ with $\Sigma = \begin{pmatrix} 4 & 3 \\ 3 & 9 \end{pmatrix}$.
    \begin{enumerate}[label=(\alph*)]
        \item \pts{4} Compute $\Corr(X, Y)$ and $\Var(2X - Y)$.
        \item \pts{4} Find the conditional distribution of $Y$ given $X = x$.
        \item \pts{4} Compute $\E[\Var(Y \mid X)]$ and $\Var(\E[Y \mid X])$, and check the law of total variance.
        \item \pts{4} Show that $X$ and $Y - \frac{3}{4}X$ are independent.
        \item \pts{4} Find a lower triangular matrix $L$ with positive diagonal entries and $LL^\top = \Sigma$, and use it to explain how to sample $(X, Y)$ from two independent standard Gaussian random variables.
    \end{enumerate}
\end{problem}

\begin{problem}[20 points]\label{pr:fin-estimation}
    Let $X_1, \ldots, X_N$ be i.i.d.\ Exponential$(\lambda)$, with the density $p(x \mid \lambda) = \lambda e^{-\lambda x}$ for $x \geq 0$, where $\lambda > 0$ is unknown.
    \begin{enumerate}[label=(\alph*)]
        \item \pts{5} Find the maximum likelihood estimator $\hat\lambda$.
        \item \pts{5} Compute the Fisher information $I(\lambda)$ of one observation, and the Cram\'er--Rao lower bound for the variance of unbiased estimators of $\lambda$.
        \item \pts{5} Use the central limit theorem and the delta method to find the limit distribution of $\sqrt{N}(\hat\lambda - \lambda)$. Compare its variance with part (b).
        \item \pts{5} Let the prior of $\lambda$ be the Gamma distribution with the shape $a$ and the rate $b$, with the density proportional to $\lambda^{a-1}e^{-b\lambda}$, the mean $\frac{a}{b}$ and, for $a \geq 1$, the mode $\frac{a - 1}{b}$. Find the posterior, its mean and the MAP estimate. Evaluate them and $\hat\lambda$ for $a = 2$, $b = 1$, $N = 10$ and $\sum_i x_i = 4$, and explain the order of the three numbers.
    \end{enumerate}
\end{problem}

\begin{problem}[20 points]\label{pr:fin-concentration-kl}
    \begin{enumerate}[label=(\alph*)]
        \item \pts{11} To estimate the click-through rate $p$ of an ad, we show it to $N$ users, who click independently with probability $p$, and use the fraction $\hat p$ of clicks. We want $P(|\hat p - p| \geq 0.05) \leq 0.05$ for every $p$. Find a sufficient $N$ (i) from Chebyshev's inequality, and (ii) from Hoeffding's inequality. (iii) Which $N$ does the normal approximation of the central limit theorem suggest, and why is it not a guarantee? (Use $\log 40 \approx 3.689$ and $1.96^2 \approx 3.84$, where $P(|Z| \geq 1.96) \approx 0.05$ for $Z \sim \mathcal{N}(0, 1)$.)
        \item \pts{9} Let $p = \mathcal{N}(0, 1)$ and $q = \mathcal{N}(1, 4)$. Compute $\KL(p \,\|\, q)$ and $\KL(q \,\|\, p)$ (use $\log 2 \approx 0.693$). Which one is larger? Explain the difference in terms of where $p$ and $q$ put their probability mass.
    \end{enumerate}
\end{problem}

\begin{problem}[Bonus, 10 points]\label{pr:fin-kkt}
    \begin{enumerate}[label=(\alph*)]
        \item \pts{4} Minimize $f(x_1, x_2) = (x_1 - 3)^2 + (x_2 - 1)^2$ subject to $x_1 + x_2 \leq 1$ and $x_2 \geq 0$. Write the constraints as $g_1 \leq 0$, $g_2 \leq 0$, state the KKT conditions, and find the minimum together with its multipliers $\mu_1, \mu_2$.
        \item \pts{1} Explain why the point found in (a) is the global minimum.
        \item \pts{5} Let $\mathbf{a} \in \mathbb{R}^n$ be nonzero. Consider the problem: minimize $\frac{1}{2}\norm{\mathbf{x}}^2$ subject to $\mathbf{a}^\top\mathbf{x} \geq 1$. Compute its dual function $d(\mu)$ for $\mu \geq 0$, solve the dual problem, and check that $d^* = p^*$ by comparing with the distance from $\mathbf{0}$ to the hyperplane $\mathbf{a}^\top\mathbf{x} = 1$.
    \end{enumerate}
\end{problem}

\end{document}
